\chapter{Introduction}\label{ch:intro}

\section{Background and Motivation}

Scientific output grows faster than any individual can read. A student or researcher who wants to know which
optimiser a paper used, how two methods compare on a benchmark, or what limitations the authors acknowledge must
locate the relevant papers, open each PDF and search its text by hand. Large language models (LLMs) can answer
such questions fluently, but a model answering from its parameters gives no indication of where an answer came
from and is known to produce plausible but unsupported statements \cite{ji2023hallucination}. For scientific use,
an answer is only useful if the reader can check it against the source.

Retrieval-augmented generation (RAG) addresses this by first retrieving passages from a document collection and
then asking the model to answer from those passages \cite{lewis2020rag,gao2023ragsurvey}. Retrieval alone does
not make an answer trustworthy: the model can still cite the wrong passage, cite a passage it was never given, or
answer when the passages do not contain the answer. This project therefore treats \emph{citation integrity} and
\emph{abstention} as engineering requirements that the server enforces and that the team measures, rather than as
properties assumed of the model.

\section{Problem Statement}

Build a software system that lets a user ask natural-language questions over a corpus of scientific papers and
returns answers in which every factual claim is traceable to a specific passage and page of a specific paper, that
explicitly declines to answer when the indexed papers do not support an answer, and whose retrieval and answer
quality are measured reproducibly. The corpus must be populated both from arXiv and from user-supplied PDF files.

\section{Objectives}

\begin{enumerate}[label=O\arabic*.]
  \item Ingest papers from the arXiv API and from uploaded PDF files, with validation and job tracking.
  \item Extract text page by page and preserve provenance from paper to page, character span and chunk.
  \item Provide semantic search over the corpus that returns ranked passages with their pages.
  \item Answer questions with a \emph{local} language model, cite each claim, verify every citation on the server,
        and abstain explicitly when evidence is insufficient.
  \item Provide single-paper summaries, cross-paper synthesis, structured extraction and comparison.
  \item Measure retrieval quality (Recall@$k$, MRR), citation validity, abstention and latency on a
        version-controlled evaluation set, and gate retrieval regressions in continuous integration.
  \item Apply software-engineering discipline throughout: requirement identifiers, traceability, architecture
        decision records, automated tests, security review and reproducible setup.
\end{enumerate}

\section{Scope and Exclusions}

The delivered system is a single-user application intended to run on one machine with Docker Compose. In scope are
the capabilities listed above together with a web interface for each of them. The following are out of scope for
the minimum viable product (MVP) and are discussed in \autoref{ch:results}: optical character recognition for
scanned PDFs, user accounts and authorisation, public deployment, and any hosted (cloud) language model. The
original proposal named a hosted model (Gemini 2.5 Flash) and a cloud deployment; the team replaced the former with
local inference through Ollama and kept deployment local (decision record ADR-0001), so that no paper or question
leaves the user's machine.

\section{Intended Users}

The primary users are students and researchers who maintain a working set of papers and want to query it. A
secondary audience is the course evaluators and any future maintainers at the institute, for whom the repository
provides requirements, decision records, test suites and operating instructions.

\section{Contributions}

The contributions of this project are engineering contributions, not new retrieval or generation methods:
\begin{itemize}
  \item A working, documented end-to-end system combining established techniques: dense retrieval
        \cite{karpukhin2020dpr,reimers2019sbert}, BM25-style lexical retrieval \cite{robertson2009bm25}, reciprocal
        rank fusion \cite{cormack2009rrf}, cross-encoder re-ranking \cite{nogueira2019rerank} and local generation.
  \item A server-side citation contract: the model may only cite passage labels; labels are resolved against the
        passages actually supplied, and an answer without a valid citation is withheld.
  \item An evaluation harness with reproducible run records, a continuous-integration retrieval gate, and audits
        of retrieval quality, security and reproducibility whose findings are reported in this document, including
        the unfavourable ones.
\end{itemize}

\section{Organisation of the Report}

\autoref{ch:litreview} reviews the techniques the system builds on. \autoref{ch:requirements} states the
requirements and use cases. \autoref{ch:architecture} presents the architecture and \autoref{ch:uml} the five UML
diagrams. \autoref{ch:implementation} describes the implementation and \autoref{ch:testing} the tests and
evaluation results. \autoref{ch:security} covers security, privacy and reliability, \autoref{ch:results} the
results, limitations and future work, and \autoref{ch:conclusion} concludes. Appendices give setup instructions,
API and configuration references, the requirements traceability matrix and test records. The report is typeset
with the TIET report class \cite{tietreport}.
